\documentclass[11pt,letterpaper]{article}
\usepackage[margin=0.75in]{geometry}
\usepackage[utf8]{inputenc}
\usepackage[T1]{fontenc}
\usepackage{lmodern}
\usepackage{helvet}
\usepackage{textcomp}
\usepackage{parskip}
\usepackage{longtable}
\usepackage{booktabs}
\usepackage{array}
\usepackage{calc}
\usepackage{amssymb}
\usepackage{enumitem}
\usepackage{titlesec}
\usepackage{xcolor}
\usepackage[colorlinks=true,urlcolor=blue,linkcolor=black]{hyperref}
\renewcommand{\familydefault}{\sfdefault}
\setlist[itemize]{leftmargin=1.3em,itemsep=1pt,topsep=2pt}
\setlist[enumerate]{leftmargin=1.5em,itemsep=1pt,topsep=2pt}
\titleformat{\section}{\large\bfseries}{}{0em}{}[\vspace{0.1em}\hrule\vspace{0.2em}]
\titlespacing*{\subsection}{0pt}{0.7em}{0.3em}
\newcolumntype{L}[1]{>{\raggedright\arraybackslash}p{#1}}
\providecommand{\tightlist}{\setlength{\itemsep}{0pt}\setlength{\parskip}{0pt}}
\setlength{\parindent}{0pt}\setlength{\parskip}{4pt}
\begin{document}
\section*{Research Statement --- \'{A}ngel Luis Robles-Fern\'andez}
\textbf{Oregon State University --- Assistant Professor, Integrative Biology}

\textbf{Environmental tolerance and persistence under variable climates.}
My research asks how evolutionary history and environmental variation shape population persistence and species interactions. At Oregon State, a central question will be: \emph{how do environmental-performance relationships and climate variability jointly determine population persistence?} I develop mathematical models, machine-learning methods, and research software that connect evidence across biological scales. My contribution will be to integrate experimental information generated by collaborators with species-occurrence records and environmental time series, linking environmental tolerance with stochastic demography. This combines an independent quantitative research agenda with the capacity to coordinate interdisciplinary biological research.

\subsection*{Scientific foundation}
In \emph{PNAS} (Robles-Fern\'andez et al., 2022), I developed machine-learning frameworks for predicting wildlife host--pathogen associations from phylogenetic, environmental, and geographic information. Their contributions differed among coronaviruses, West Nile virus, and avian malaria, demonstrating why interaction prediction must account for each system's biology. Collaborations extended this approach to ectoparasites and parasitic plants. In \emph{Ecography} (Hern\'andez, Robles-Fern\'andez \& Upham, 2025; co-first authors), we developed the Suitability Prevalence Area method, connecting Late Quaternary environmental stability with present-day population genetic diversity in \emph{Sciurus aberti}.

In Daniel Reuman's laboratory at the University of Kansas, I co-developed XSDM, which incorporates stochastic demography and interannual environmental variability into species distribution modeling (Berti et al., \emph{bioRxiv} preprint). In our analyses, accounting for variability reduced potential ranges by 22\% on average and up to 45\% for some species. I created and maintain the \texttt{xsdm} R package and develop its Bayesian counterpart, \texttt{xsdmbayes}. These frameworks connect growth--environment relationships with long-term stochastic growth, providing a basis for incorporating experimentally characterized environmental responses into distribution forecasts.

My training in physics and evolutionary biology and experience in industrial machine learning connect biological questions with algorithms and research software. My LACS literature-classification project received second prize in the 2022 GBIF Ebbe Nielsen Challenge, establishing a foundation for AI-assisted data acquisition.

\subsection*{Research agenda at Oregon State}
\textbf{1. Connecting experimental performance with demographic predictions.}
In an ongoing collaboration, we propose translating experimental performance--environment curves into priors for Bayesian XSDM. This work has not yet been published. The existing framework accommodates configurable priors, allowing comparison of models using reference priors without experimental information against models using performance-informed priors. My role is to develop the statistical integration and software alongside collaborators who generate and interpret experimental evidence. The mapping from experimental endpoints to model parameters must make units, temporal scales, and the connection to population growth explicit.

At OSU, I will test whether experimentally characterized response breadth and asymmetry explain differences in sensitivity to environmental variability, using thermal responses as an example. Comparisons will evaluate held-out geographic predictions, posterior uncertainty, and sensitivity to prior choices. Rather than assuming experiments necessarily improve forecasts, I will ask when their information transfers to field conditions and when discrepancies reveal limitations of the model or the experimental-to-demographic mapping. This provides a concrete bridge between organismal performance, environmental tolerance, and ecological persistence.

\textbf{2. AI-enabled software and databases for integrative biology.}
Which combinations of experimental, genetic, environmental, and interaction evidence support transferable comparative models? I will extend my current software stack to build research-ready datasets for this question. \href{https://ecoseek.org}{EcoSeek} connects ecological tools and scientific agents; \href{https://genominer.ecoseek.org}{Genominer} organizes candidate population-genetic datasets; \href{https://xbioclim.org}{xbioclim} provides year-resolved climate layers. AI-assisted extraction will be assessed against expert-curated records, with source provenance and taxonomic harmonization. These resources will support performance-informed analyses and comparative tests of relationships such as environmental stability and genetic diversity. Combining interaction models with XSDM will also allow me to examine changes in host and vector persistence, distinguishing opportunities for encounters from evidence of infection or host competence.

Parallel methodological developments include \texttt{arfun}, with Jorge Sober\'on, for ancestral reconstruction of environmental niche position and shape, and \texttt{ancemb} for Bayesian reconstruction of ancestral protein embeddings. Both remain in development and require independent biological validation; they extend the evolutionary methods available to the program rather than constitute completed findings about organismal function.

\textbf{3. Evaluating AI-driven scientific workflows.}
Through EcoSeek, the DiDAL critique-and-revision protocol, and \href{https://litreview.org}{litreview}, I am developing scientist-led evidence-synthesis workflows. \href{https://sciev.org}{Sciev}, my project on structured scientific decision models, adds a route to candidate selection and evidence assessment. I will evaluate these tools on biological tasks using expert-labeled references, held-out data, and controls with missing or conflicting evidence. The goal is to determine whether AI makes research more reproducible and efficient, while scientists retain responsibility for hypotheses, interpretation, and conclusions.

\subsection*{Collaboration, funding, and training}
My program would complement organismal expertise in Integrative Biology through jointly defined biological questions and experimental interpretation, with my laboratory leading statistical modeling, data integration, and reusable software. I will also seek transdisciplinary partnerships with conservation and disease-data practitioners to define useful data products and evaluation criteria.

I will pursue external funding for integrative ecological and evolutionary modeling, biological data infrastructure, and AI-enabled research methods. Initial proposals will build on existing software, published studies, and collaborative work connecting performance information with demographic forecasts. Undergraduates will enter through supported data and analysis projects; graduate students will develop independent questions, methods, and collaborations. Documented software, reproducible datasets, and cross-disciplinary training will sustain a laboratory that links detailed biological evidence with broader ecological and evolutionary inference.
\end{document}
